\documentclass[a4paper,11pt]{article}

\usepackage[utf8]{inputenc}
\usepackage[T1]{fontenc}
\usepackage[spanish]{babel}
\usepackage[margin=2.5cm]{geometry}
\usepackage{amsmath,amssymb,amsthm}
\usepackage[shortlabels]{enumitem}
\usepackage{parskip}
\usepackage{xcolor}
\usepackage{listings}
\usepackage{booktabs}

% ─── estilo de listados (igual que NotasUltimaClase.tex) ────────────────────
\lstset{
  language=Python,
  basicstyle=\ttfamily\small,
  keywordstyle=\color{blue!70!black}\bfseries,
  commentstyle=\color{gray}\itshape,
  stringstyle=\color{red!60!black},
  showstringspaces=false,
  breaklines=true,
  frame=single,
  framerule=0.4pt,
  numbers=left,
  numbersep=8pt,
  numberstyle=\tiny\color{gray},
  xleftmargin=15pt
}

\theoremstyle{remark}
\newtheorem*{nota}{Nota}

% ────────────────────────────────────────────────────────────────────────────
\begin{document}

\section*{Solución --- Parte A}

Se identifica la correspondencia entre las variables del modelo general
(notas de clase) y las del problema:
\begin{align*}
  y_i &\equiv M_i \quad \text{(magnitud absoluta en banda }V\text{)},\\
  x_i &\equiv \log_{10}(P_i) \quad \text{(logaritmo del período en días)}.
\end{align*}
El parámetro de error $\sigma = 0.3$ mag es conocido; las distribuciones
previas son $a,b\sim\mathcal{N}(0,10^2)$ independientes, con
$\sigma_a = \sigma_b = 10$.

% =====================================================
\subsection*{Preparación de los datos}
% =====================================================

El archivo CSV de formato CDS (VizieR) contiene dos filas de encabezado
(unidades y separadores de columnas) que quedan como texto.
Convirtiendo cada columna relevante con \texttt{errors='coerce'},
esas filas se transforman en \texttt{NaN} y se eliminan junto con
cualquier observación incompleta.

\begin{nota}
  A pesar de la etiqueta \texttt{mas}, los valores de \texttt{plx}
  están almacenados en microssegundos de arco ($\mu$as). Se dividen por
  $10^3$ para obtener el parálaje en mas antes de aplicar el módulo de
  distancia $\mu = 10 - 5\log_{10}(\varpi[\text{mas}])$, lo que da
  $M_V = V + 5\log_{10}(\varpi) - 10$.
\end{nota}

\begin{lstlisting}
import numpy as np, pandas as pd, matplotlib.pyplot as plt

df = pd.read_csv('mw_cepheids.csv')
for col in ['Per', 'plx', 'Vmag']:
    df[col] = pd.to_numeric(df[col], errors='coerce')  # texto -> NaN
df.dropna(subset=['Per', 'plx', 'Vmag'], inplace=True) # elimina NaN

plx_mas = df['plx'].values / 1000                         # paralaje [arcsec]
x  = np.log10(df['Per'].values)                           # log10(P_i) [adim]
y  = df['Vmag'].values + 5*np.log10(plx_mas) + 5        # M_i [mag]
n  = len(x)
sigma = 0.3                                                # desv. est. conocida
\end{lstlisting}

% =====================================================
\begin{enumerate}[a)]

% ─────────────── a ───────────────
\item \textbf{Gráfico de dispersión.}

\begin{lstlisting}
plt.figure()
plt.scatter(x, y, s=10, alpha=0.6)
plt.xlabel(r'$\log_{10}(P)$ [adim.]')
plt.ylabel(r'$M_V$ [mag]')
plt.gca().invert_yaxis()   # magnitudes: escala invertida
plt.tight_layout(); plt.savefig('dispersion.pdf')
\end{lstlisting}

\textbf{Tendencia observada:} se aprecia una relación lineal negativa
(Ley de Leavitt): a mayor período, menor magnitud absoluta (estrella más
luminosa). La dispersión alrededor de la tendencia es pequeña, lo que
justifica el modelo de regresión lineal propuesto.

% ─────────────── b ───────────────
\item \textbf{Función de verosimilitud.}

Como $y_i\mid M(a,b)\sim\mathcal{N}(a+bx_i,\,\sigma^2)$ de forma
independiente para cada observación, la verosimilitud conjunta es
$L(y\mid M(a,b))=\prod_i f(y_i\mid M(a,b))$. Descartando la constante
de normalización:
\begin{equation}
  L\!\left(y\mid M(a,b)\right)
  \;\propto\;
  \exp\!\left(-\frac{1}{2\sigma^2}\sum_{i=1}^{n}
              \bigl(y_i - a - bx_i\bigr)^2\right).
  \label{eq:vero}
\end{equation}

% ─────────────── c ───────────────
\item \textbf{Distribución posterior conjunta.}

Por el teorema de Bayes, $P(M(a,b)\mid y)\propto
L(y\mid M(a,b))\cdot f(a,b)$, con prior
$f(a,b)\propto\exp\!\bigl(-\tfrac{a^2+b^2}{200}\bigr)$:
\begin{equation}
  P\!\left(M(a,b)\mid y\right)
  \;\propto\;
  \exp\!\left(
    -\frac{1}{2\sigma^2}\sum_{i}(y_i-a-bx_i)^2
    -\frac{a^2+b^2}{200}
  \right).
  \label{eq:post}
\end{equation}

% ─────────────── d ───────────────
\item \textbf{Estimación MAP.}

Se toma el logaritmo de~\eqref{eq:post} y se igualan a cero las
derivadas parciales respecto de $a$ y $b$.
Las ecuaciones resultantes son \textbf{lineales} en $(\hat{a},\hat{b})$
y forman el sistema $\mathbf{M}\vec\theta = \vec Z$:
\begin{equation}
  \underbrace{\begin{bmatrix}
    n + \dfrac{\sigma^2}{100} & \displaystyle\sum_i x_i \\[8pt]
    \displaystyle\sum_i x_i   & \displaystyle\sum_i x_i^2
                                +\dfrac{\sigma^2}{100}
  \end{bmatrix}}_{\mathbf{M}}
  \underbrace{\begin{bmatrix}\hat{a}\\\hat{b}\end{bmatrix}}_{\vec\theta}
  =
  \underbrace{\begin{bmatrix}
    \displaystyle\sum_i y_i \\[4pt]
    \displaystyle\sum_i x_iy_i
  \end{bmatrix}}_{\vec Z}.
  \label{eq:sistema}
\end{equation}

Los términos $\sigma^2/100$ en la diagonal provienen de derivar el
término de la previa $-\tfrac{a^2+b^2}{200}$ e incorporan regularización
tipo \emph{ridge} con parámetro efectivo $\lambda = \sigma^2/100 = 0.0009$.

\begin{lstlisting}
sum_x  = x.sum();    sum_x2 = (x**2).sum()   # estadisticos suficientes
sum_y  = y.sum();    sum_xy = (x*y).sum()

M_mat = np.array([[n + sigma**2/100, sum_x ],
                   [sum_x,           sum_x2 + sigma**2/100]])  # matriz M
Z_vec = np.array([sum_y, sum_xy])                              # vector Z

a_map, b_map = np.linalg.solve(M_mat, Z_vec)   # resuelve M theta = Z
\end{lstlisting}

% ─────────────── e ───────────────
\item \textbf{La posterior es una normal bivariada.}

Se expande el cuadrado en~\eqref{eq:post}:
\[
  \sum_i(y_i-a-bx_i)^2
  = na^2 + 2ab\!\sum_i x_i + b^2\!\sum_i x_i^2
    - 2a\!\sum_i y_i - 2b\!\sum_i x_iy_i
    + \underbrace{\sum_i y_i^2}_{\text{cte en }(a,b)}.
\]
Sustituyendo en~\eqref{eq:post} y agrupando con el término de la previa,
el exponente toma la forma cuadrática:
\begin{equation}
  \ln P(M(a,b)\mid y)
  = -\frac{1}{2}\!\left(Aa^2 + 2Bab + Cb^2 - 2Da - 2Eb\right)
    + \mathrm{cte},
  \label{eq:forma-cuad}
\end{equation}
con coeficientes
\begin{align}
  A &= \frac{n}{\sigma^2}+\frac{1}{100}, &
  B &= \frac{\sum_i x_i}{\sigma^2}, &
  C &= \frac{\sum_i x_i^2}{\sigma^2}+\frac{1}{100},
  \label{eq:ABC}\\
  D &= \frac{\sum_i y_i}{\sigma^2}, &
  E &= \frac{\sum_i x_iy_i}{\sigma^2}.&&
  \label{eq:DE}
\end{align}

La expresión~\eqref{eq:forma-cuad} es exactamente el núcleo de una
distribución normal bivariada (cf.\ notas de clase).  
Por tanto:
\begin{equation}
  (a,b)\mid y \;\sim\; \mathcal{N}_2\!\bigl((\hat{a},\hat{b}),\,\Sigma\bigr),
  \label{eq:posterior-normal}
\end{equation}
donde $(\hat{a},\hat{b})$ es la solución de~\eqref{eq:sistema} (media
$=$ moda para una normal), la \textbf{matriz de precisión} es
\begin{equation}
  \Sigma^{-1}
  = \begin{bmatrix}A & B \\ B & C\end{bmatrix},
  \qquad
  \Sigma
  = \frac{1}{AC-B^2}\begin{bmatrix}C&-B\\-B&A\end{bmatrix}.
  \label{eq:cov}
\end{equation}

\begin{lstlisting}
A = n/sigma**2 + 1/100
B = sum_x/sigma**2
C = sum_x2/sigma**2 + 1/100

Prec  = np.array([[A, B], [B, C]])  # matriz de precision Sigma^{-1}
Sigma = np.linalg.inv(Prec)         # covarianza posterior
\end{lstlisting}

% ─────────────── f ───────────────
\item \textbf{Muestras de la distribución posterior.}

Se usa directamente el resultado~\eqref{eq:posterior-normal}:
\begin{lstlisting}
np.random.seed(42)
theta_s = np.random.multivariate_normal([a_map, b_map], Sigma, size=10_000)
a_s, b_s = theta_s[:, 0], theta_s[:, 1]   # muestras marginales de a y b
\end{lstlisting}

% ─────────────── g ───────────────
\item \textbf{Intervalos de credibilidad del 95\% para $a$ y $b$.}

Se obtienen como los percentiles 2.5 y 97.5 de las muestras:
\begin{lstlisting}
ci_a = np.percentile(a_s, [2.5, 97.5])   # IC 95% para a
ci_b = np.percentile(b_s, [2.5, 97.5])   # IC 95% para b
\end{lstlisting}

% ─────────────── h ───────────────
\item \textbf{Distribución posterior predictiva para $P = 12$ días.}

Para cada muestra $(\tilde{a},\tilde{b})$ de la posterior se predice
$\tilde{M} = \tilde{a} + \tilde{b}\log_{10}(12)$; la colección de
$\tilde{M}$ aproxima la distribución posterior de la magnitud absoluta:
\begin{lstlisting}
x_new  = np.log10(12)                          # log10(P) para P = 12 dias
M_pred = a_s + b_s * x_new                    # distribucion predictiva posterior
mu_pred, ci_pred = M_pred.mean(), np.percentile(M_pred, [2.5, 97.5])
\end{lstlisting}

% ─────────────── i ───────────────
\item \textbf{Relación período--luminosidad con banda de credibilidad.}

Se evalúa cada muestra $(\tilde{a},\tilde{b})$ en una grilla de
$\log_{10}(P)$ y se toman los percentiles 2.5 y 97.5 columna a columna:
\begin{lstlisting}
x_grid = np.linspace(x.min(), x.max(), 300)
M_grid = a_s[:, None] + b_s[:, None]*x_grid  # (10000, 300): curva por muestra
lo, hi = np.percentile(M_grid, [2.5, 97.5], axis=0)

plt.figure()
plt.scatter(x, y, s=10, alpha=0.5, label='Datos')
plt.plot(x_grid, a_map + b_map*x_grid, 'r-', lw=1.5, label='MAP')
plt.fill_between(x_grid, lo, hi, alpha=0.3, label='Banda 95\%')
plt.xlabel(r'$\log_{10}(P)$'); plt.ylabel(r'$M_V$ [mag]')
plt.gca().invert_yaxis(); plt.legend(); plt.tight_layout()
plt.savefig('pl_band.pdf')
\end{lstlisting}

% ─────────────── j ───────────────
\item \textbf{Comparación con regresión por mínimos cuadrados (OLS).}

Las ecuaciones normales de OLS son idénticas a~\eqref{eq:sistema} salvo
por los términos de regularización diagonal:
\[
  \underbrace{\begin{bmatrix}n & \sum x_i \\ \sum x_i & \sum x_i^2
  \end{bmatrix}}_{\mathbf{M}_{\text{OLS}}}
  \begin{bmatrix}\hat{a}_{\text{OLS}}\\\hat{b}_{\text{OLS}}\end{bmatrix}
  = \begin{bmatrix}\sum y_i\\\sum x_iy_i\end{bmatrix}.
\]

\begin{lstlisting}
X_mat = np.column_stack([np.ones(n), x])                    # diseno [1 | x]
a_ols, b_ols = np.linalg.lstsq(X_mat, y, rcond=None)[0]    # estimacion OLS
\end{lstlisting}

\textbf{Semejanzas.} Como $\sigma^2/100 = 0.0009 \ll n$, la
perturbación diagonal es despreciable y los estimadores MAP y OLS son
numéricamente casi idénticos.

\textbf{Diferencias.}
\begin{itemize}
  \item El enfoque bayesiano produce la distribución posterior completa
    \eqref{eq:posterior-normal}, por lo que cualquier probabilidad sobre
    $(a,b)$ o sobre una predicción futura se calcula directamente como
    área bajo dicha densidad (intervalos de \emph{credibilidad}).
  \item OLS produce intervalos de \emph{confianza} frecuentistas: la
    interpretación es que, en repetidas muestras, el $95\%$ de los
    intervalos construidos cubrirían el verdadero parámetro; $(a,b)$ no
    es una variable aleatoria en ese marco.
  \item La incertidumbre en $(a,b)$ se propaga de forma natural a
    cualquier cantidad de interés (inciso~h) mediante las muestras
    posteriores, sin necesidad de aproximaciones adicionales como el
    método delta o remuestreo.
\end{itemize}

\end{enumerate}

\end{document}